\frontchapter{Foreword and Prologue}

\ssection{The Dedication}
Shelley dedicates the book to his church history students, ``who pressed the question of significance''---the question the whole book tries to answer.

\ssection{Foreword, by R. L. Hatchett}
Hatchett, professor of theology and philosophy at Houston Baptist University, revised the book for this
edition. He states what the revision is for: evangelical readers should come to know their own history and, through it, the wider Christian world. He credits its success to Shelley himself, to his character and concern for the
church, his training and eye as a historian, his vivid prose and his long classroom experience, and
to his discipline. Shelley wrote for a popular audience in plain language, chose representative
stories instead of aiming at completeness, made each story readable on its own, and kept citations out
of the reader's way: references are in the endnotes, and wider debts in the reading lists that close
each chapter.

Hatchett kept this informal style. He added endnotes and marked the books he newly recommends with
asterisks. Apart from small changes throughout, he added material on Gnosticism and its continuing
relevance, on the theology of the early church and of the Reformation, and above all on the rapid
global spread and transformation of Christianity since 1900, which fills three new chapters. Recent
history is hard to write, he says, but the surprising situation of Christianity in the West and in the
Global South needs explaining as much as retelling.

He adds a word on names and numbers. Terms such as ``Global South'' and ``majority world'' are
unsatisfactory, but he uses the common ones for want of better. The ``West'' means North America,
England and Western Europe; the ``Global South'' means the Third World---South America, Africa and
Asia, including places north of the equator such as China. ``Far East'' and ``Middle East'' betray the
same Western point of view. His point is that evangelicals in the United States must widen their horizon
if they are to share in the work now flourishing across the world. He usually cites modest figures from
established sources and believes them too low: watching his late friend Walter Lumpkin try to count the
non-denominational churches of Houston taught him that the mainline churches count well, while the most
vital groups are the ones that are missed.

He thanks the students who helped with typing and research, several friends with wide experience of the world, the editors at Thomas Nelson who recruited him and saw the project through, and most of all his wife Debbie.

\ssection{Prologue, by Bruce L. Shelley}
Shelley begins with a \emph{Peanuts} cartoon that hung on his study door. Sally, writing a school theme
on church history, starts ``at the very beginning'': ``Our pastor was born in 1930.'' Many Christians,
he says, suffer from the same historical amnesia; for them the time between the apostles and today is
a blank. God, however, cares about time: Israel was told to explain the Passover to its children and to
keep a jar of manna for the generations to come.

Ignorance of Christian history has three costs. Believers fall for cults because they take a distortion
for the real thing. Others, with nothing to compare themselves to, become proud and defend their own
way and party as the best. And many serve in ministry without a wider context for judging how to use
their time and effort. One survey will not cure all this, but history helps to separate the transient
from the permanent and fads from basics.

The book is written for laypeople. Forty years in the seminary classroom have convinced Shelley that the college graduate who begins theological study and the lay reader with a job in engineering or sales, who gets through a handful of books a year, form a single audience. He has borrowed freely from other writers with one aim, to keep the story
moving; the most helpful books are listed after each chapter and the major quotations at the end.

Because clarity is the first law of learning, the story is divided into ages, each marked by a change
in the conditions of the church's life, though every age carries remains of the past and seeds of the
future. Whoever wants the plot will find it in the short introductions that open the nine ages. Each
chapter treats a single issue, put as a question after an opening anecdote, so that it can be read
alone like an encyclopedia article. This leaves gaps, and some important people and events are missing,
but it shows the present-day significance of the past: many of today's issues are not new. There is
much biography because readers like to meet people, so ideas are wrapped in personalities.

Is the church a movement or an institution? Shelley answers: both. He therefore tells of missionary
expansion as well as papal politics and does not hold to a strict definition of ``church.'' God's people are always pulled two ways, toward the ideal of one universal fellowship and toward the actual congregation that exists at a given place and date. Mission requires institutions (rules, officers, buildings); when these get in the way of the gospel, renewal movements appear and call the church back to its task. The book shows how often that has happened.

He closes with thanks to the helpers of the third edition: Daniel Hallock for research on the Internet,
his colleague David Buschart for many titles in the reading lists and for conversation about the
emerging churches, and his colleague and successor Scott Wenig for further titles, discussion and that
edition's foreword.
